\section*{Задачи}
\addcontentsline{toc}{section}{Задачи}

\begin{enumerate}
    \item \textit{(Демо 2026, №4)} В базисе $e_1 = \binom{3}{1}, e_2 = \binom{2}{-1}$ билинейная форма $B(x,y)$ имеет матрицу 
    $B = \begin{pmatrix} -1 & 1 \\ -3 & 4 \end{pmatrix}$. 
    Найти матрицу билинейной формы $B(x,y)$ в базисе $\hat{e}_1 = \binom{4}{3}, \hat{e}_2 = \binom{1}{1}$.

    \item \textit{(Демо 2026, №5)} В базисе $e_1 = \binom{3}{1}, e_2 = \binom{2}{-1}$ квадратичная форма $Q(x)$ имеет матрицу 
    $B = \begin{pmatrix} -1 & 1 \\ -1 & 4 \end{pmatrix}$. 
    Найти матрицу квадратичной формы $Q(x)$ в базисе $\hat{e}_1 = \binom{4}{3}, \hat{e}_2 = \binom{1}{1}$.

    \item \textit{(Проскуряков, №1244)} Найти симметричную билинейную форму, соответствующую квадратичной форме: 
    $x_1^2 + 2x_1 x_2 + 2x_2^2 + 4x_2 x_3 + 5x_3^2$.

    \item \textit{(Проскуряков, №1249)} Найти матрицу следующей квадратичной формы: 
    $x_1^2 - 2x_1 x_2 + x_2^2$.

    \item \textit{(Демо 2026, №6)} Привести квадратичную форму 
    $2x_1 x_2 - 2x_1 x_3 - x_1^2 + x_2^2 + 5x_3^2$ 
    к нормальному виду с помощью метода Лагранжа. Определить ранг и индексы инерции. Указать соответствующее линейное преобразование.

    \item \textit{(Проскуряков, №1571)} Пользуясь методом Лагранжа, привести квадратичную форму к каноническому виду: 
    $x_1^2 + x_2^2 + 3x_3^2 + 4x_1 x_2 + 2x_1 x_3 + 2x_2 x_3$.

    \item \textit{(Проскуряков, №1574)} Пользуясь методом Лагранжа, привести квадратичную форму к каноническому виду: 
    $x_1 x_2 + x_1 x_3 + x_2 x_3$.

    \item \textit{(Проскуряков, №1586)} Найти нормальный вид следующей квадратичной формы в вещественной области: 
    $x_1^2 + x_2^2 + 3x_3^2 + 4x_1 x_2 + 2x_1 x_3 + 2x_2 x_3$.

    \item \textit{(Демо 2026, №7 и №18)} Привести квадратичную форму 
    $k = x_1^2 - 6x_1x_2 - 2x_1x_3 + x_2^2 + 2x_2x_3 + 5x_3^2$ 
    к каноническому виду посредством ортогональной замены координат. Определить ранг и индексы инерции. Указать соответствующее линейное преобразование.

    \item \textit{(Проскуряков, №1370)} Найти канонический вид, к которому приводится следующая квадратичная форма посредством ортогонального преобразования: 
    $x_1^2 + x_2^2 + x_3^2 + 4x_1 x_2 + 4x_1 x_3 + 4x_2 x_3$.

    \item \textit{(Проскуряков, №1374)} Найти ортогональное преобразование, приводящее квадратичную форму к каноническому виду (и выписать этот канонический вид): 
    $2x_1^2 + x_2^2 - 4x_1 x_2 - 4x_2 x_3$.

    \item \textit{(Демо 2026, №8)} Исследовать квадратичную форму на положительную или отрицательную определенность в зависимости от параметра: 
    $k = (\alpha - 1)x_1^2 + (2\alpha - 2)x_1x_2 - 2\alpha x_1x_3 + 2\alpha x_2^2 - 2\alpha x_2x_3 + (\alpha - 2)x_3^2$.

    \item \textit{(Проскуряков, №1598)} Найти все значения параметра $\lambda$, при которых квадратичная форма положительно определена: 
    $5x_1^2 + x_2^2 + \lambda x_3^2 + 4x_1 x_2 - 2x_1 x_3 - 2x_2 x_3$.

    \item \textit{(Проскуряков, №1600)} Найти все значения параметра $\lambda$, при которых квадратичная форма положительно определена: 
    $x_1^2 + 4x_2^2 + x_3^2 + 2\lambda x_1 x_2 + 10x_1 x_3 + 6x_2 x_3$.
\end{enumerate}
